\SPBGContentSlide
  {Строгое линейное отделение}
  {%
    \SmallBody
    Рассмотрим два конечных множества точек и метки классов:
    \[
      P_1=\{p_j\}_{j=1}^{s},
      \qquad
      P_2=\{p_j\}_{j=s+1}^{m},
      \qquad
      \xi_j=\begin{cases}
        +1,& j=1,\ldots,s,\\
        -1,& j=s+1,\ldots,m.
      \end{cases}
    \]

    Гиперплоскость задаётся уравнением
    \[
      \langle w,p\rangle+\beta=0,
      \qquad
      w\in\mathbb{R}^{n},\quad w\neq0.
    \]

    Множества \(P_1\) и \(P_2\) допускают строгое линейное отделение,
    если существуют \(w\neq0\) и \(\beta\), такие что
    \[
      \begin{aligned}
        \langle w,p_j\rangle+\beta&>0,
        &&j=1,\ldots,s,\\
        \langle w,p_j\rangle+\beta&<0,
        &&j=s+1,\ldots,m.
      \end{aligned}
    \]

  }
